%---------------------------------------------------------------------------%
%->> Frontmatter
%---------------------------------------------------------------------------%
%-
%-> 生成封面
%-

\maketitle% 生成中文封面
\MAKETITLE% 生成英文封面
%-
%-> 作者声明
%-
\makedeclaration% 生成声明页
%-
%-> 中文摘要
%-
\intobmk\chapter*{摘\quad 要}% 显示在书签但不显示在目录
\setcounter{page}{1}% 开始页码
\pagenumbering{Roman}% 页码符号

<<<<<<< HEAD
多光子荧光寿命成像技术结合了多光子显微成像的深层组织穿透能力、内禀三维光学切片能力，以及荧光寿命成像对探针浓度、激发功率和组织衰减不敏感的定量优势，在活体神经成像和复杂微环境分析中具有重要应用价值。围绕活体生物成像对高分辨率、高灵敏度和高时序精度的需求，本文开展了共聚焦与多光子荧光寿命成像系统的理论分析、光机系统构建、超快脉冲调控、同步采集链路开发及寿命拟合研究。

本文首先从点扫描显微成像中的像素停留时间、采样分辨率、点扩散函数及有效激发体积等问题出发，分析了扫描视场、物镜参数与成像性能之间的关系，并为后续系统设计提供了理论依据。在此基础上，构建了模块化激光扫描共聚焦荧光寿命显微镜，提出磁吸式快拆多色光学架构，实现了多波段激发与探测通道的高重复性快速切换，并建立了仪器响应函数标定、同步延迟校准及数据采集流程。

面向活体大范围与高帧率成像需求，本文进一步设计并构建了两套多光子显微平台，实现了振镜-振镜与振镜-共振镜双扫描模式的无缝切换。针对飞秒脉冲在显微系统中的色散展宽问题，设计了基于棱镜对的脉冲压缩与色散补偿装置，并结合自建自相干仪完成了物镜后脉宽表征与补偿效果验证。同时，本文打通了单路与三路同步的时间相关单光子计数采集链路，优化了硬件死区时间限制，建立了从同步采集、像素重分配到逐像素寿命拟合的完整数据处理流程，并开发了与商用软件结果一致性较好的寿命拟合算法。

本文完成了从点扫描理论分析、光机系统搭建、多光子激发与脉冲调控，到荧光寿命采集与数据处理链路建立的系统研究，为后续面向深层组织功能成像与活体神经活动记录的系统优化和应用拓展奠定了基础。
=======
多光子荧光寿命成像技术融合了多光子显微成像的深层组织穿透能力、三维光学切片能力以及荧光寿命成像对探针浓度、激发功率和组织衰减不敏感的定量优势，在活体神经成像、功能表征和复杂微环境分析中具有重要应用价值。围绕活体生物成像对高分辨率、高灵敏度和高时序精度的需求，本文系统开展了共聚焦与多光子荧光寿命成像系统的理论分析、光机架构创新、同步采集闭环和寿命拟合研究。

首先，本文从点扫描显微成像中的速度、采样分辨率、点扩散函数及有效激发体积出发，分析了扫描视场、像素停留时间、物镜参数与系统成像性能之间的关系，并对点扫描显微系统中的典型扫描策略进行了比较。在此基础上，结合多光子激发和荧光寿命成像需求，对显微系统的扫描结构与光学参数进行了设计分析，为后续系统实现提供了理论依据。

其次，针对传统共聚焦系统多色扩展性差、光路切换易失调的问题，本文完成了激光扫描共聚焦荧光寿命显微镜的预研搭建，给出了多色激发与系统表征方案，并围绕荧光寿命成像的实现需求，建立了仪器响应函数标定、同步延迟校准及数据采集流程；自主设计并搭建了模块化激光扫描共聚焦荧光寿命显微镜，创新性地提出了磁吸式快拆多色光学架构，实现了多波段激发与探测通道的高重复性快速切换。同时，结合Zemax软件完成了10倍放大0.4数值孔径大视场物镜的自主光学仿真设计，为大视场高通量扫描提供了理论验证。

再次，面向活体大范围与高帧率成像的复杂需求，本文设计并构建了两套多光子显微平台。其中，快速切换多光子显微系统采用了``整体显微镜随动''的高精度 XY 位移台架构，并自主设计中继光路，实现了振镜-振镜与振镜-共振镜双扫描模式的空间共轭与无缝切换。针对飞秒脉冲色散问题，设计了基于棱镜对的脉冲压缩与色散补偿装置，并进一步分析了多通腔压缩方案，为提升多光子激发效率提供了技术支撑。

最后，在荧光寿命数据处理与系统表征方面，本文不仅打通了单路与三路同步的时间相干单光子计数采集链路，优化了硬件死区时间限制，还开发了与商用软件精度相当的逐像素寿命拟合算法。此外，为进一步完善多光子成像的基础支撑，本文自主搭建了基于紫外响应宽禁带半导体探测器的双光子干涉自相干仪用于超快脉冲精准表征，并提出了一种双椭圆级联等光程扫描架构，有效校正了传统扫描引起的视场畸变与速率非均匀性。

本文的工作完成了从前沿扫描理论分析、创新光机系统设计、多光子激发与脉冲调控，到荧光寿命采集和自研数据处理链路的全面构建，为后续面向活体神经活动记录和多模态功能成像的系统优化与应用拓展奠定了坚实的软硬件基础。
>>>>>>> c4bb8a8e4d07c80bffd10fe9aeb672bafe56fbee

\keywords{多光子显微成像，多色共聚焦显微镜，荧光寿命成像，时间相关单光子计数，色散补偿}% 中文关键词
%-
%-> 英文摘要
%-
\intobmk\chapter*{Abstract}% 显示在书签但不显示在目录
<<<<<<< HEAD
Multiphoton fluorescence lifetime imaging combines the deep-tissue penetration and intrinsic three-dimensional optical sectioning capability of multiphoton microscopy with the quantitative advantages of fluorescence lifetime imaging (FLIM), which is largely insensitive to probe concentration, excitation power, and tissue attenuation. It is therefore valuable for in vivo neural imaging and complex microenvironment analysis. To meet the demands of in vivo biological imaging for high spatial resolution, high sensitivity, and high temporal precision, this thesis investigates the theoretical analysis, system construction, ultrafast pulse control, synchronized acquisition, and lifetime fitting of confocal and multiphoton FLIM systems.

Starting from the fundamental issues of point-scanning microscopy, including pixel dwell time, sampling resolution, point spread function, and effective excitation volume, this thesis analyzes the relationships among scanning field of view, objective parameters, and imaging performance, providing the theoretical basis for subsequent system design. On this basis, a modular laser-scanning confocal FLIM system was constructed. A magnetic quick-release multi-color optical architecture was proposed to enable highly repeatable and rapid switching of multi-band excitation and detection channels, and workflows for instrument response function calibration, synchronization-delay calibration, and data acquisition were established.

For large-area and high-frame-rate imaging, two multiphoton microscopy platforms were further designed and built, enabling seamless switching between Galvo-Galvo and Galvo-Resonant scanning modes. To compensate for femtosecond pulse broadening in the microscope, a prism-pair-based pulse-compression and dispersion-compensation scheme was implemented, and its effectiveness was verified through pulse-width characterization after the objective using a self-built autocorrelator. In addition, single-channel and three-channel synchronized time-correlated single-photon counting acquisition pipelines were established, hardware dead-time limitations were optimized, and a complete data-processing workflow from synchronized acquisition and pixel reassignment to pixel-wise lifetime fitting was developed. A lifetime fitting algorithm with good consistency with commercial software was also achieved.

Overall, this thesis completes a systematic study spanning point-scanning theory, opto-mechanical system construction, multiphoton excitation and pulse control, and fluorescence lifetime acquisition and processing, providing a foundation for future optimization and applications in deep-tissue functional imaging and in vivo neural activity recording.
=======
Multiphoton fluorescence lifetime imaging integrates the deep tissue penetration and intrinsic three-dimensional optical sectioning capabilities of multiphoton microscopy with the quantitative advantages of fluorescence lifetime imaging (FLIM), which is insensitive to probe concentration, excitation power, and tissue attenuation. It holds significant application value in in vivo neural imaging, functional characterization, and complex microenvironment analysis. Addressing the demands of in vivo biological imaging for high resolution, high sensitivity, and high temporal precision, this thesis systematically conducts research on theoretical analysis, opto-mechanical architectural innovation, synchronized acquisition closed-loops, and lifetime fitting for confocal and multiphoton FLIM systems.

First, starting from the speed, sampling resolution, point spread function (PSF), and effective excitation volume in point-scanning microscopy, this thesis analyzes the relationships among scanning field of view (FOV), pixel dwell time, objective parameters, and system imaging performance. Typical scanning strategies in point-scanning microscopy systems are also compared. Building upon this, by combining the requirements of multiphoton excitation and FLIM, the scanning structure and optical parameters of the microscopy system are designed and analyzed, providing a theoretical basis for subsequent system implementation.

Second, targeting the issues of poor multi-color scalability and optical switching misalignment in traditional confocal systems, this thesis completes the preliminary construction of a laser-scanning confocal FLIM system, providing schemes for multi-color excitation and system characterization. Centering on the implementation requirements of FLIM, workflows for instrument response function (IRF) calibration, synchronization delay calibration, and data acquisition are established. A modular laser-scanning confocal FLIM system is independently designed and constructed, innovatively proposing a magnetic quick-release multi-color optical architecture to achieve highly repeatable and rapid switching of multi-band excitation and detection channels. Meanwhile, an autonomous optical simulation design of a 10$\times$ magnification, 0.4 numerical aperture (NA) large-FOV objective is completed using Zemax, providing theoretical validation for large-FOV, high-throughput scanning.

Third, facing the complex demands of large-area and high-frame-rate in vivo imaging, this thesis designs and constructs two multiphoton microscopy platforms. Among them, the fast-switching multiphoton microscopy system employs a breakthrough "whole-microscope-moving" architecture integrated onto a high-precision XY translation stage. By independently designing the relay optical path, it achieves spatial conjugation and seamless switching between Galvo-Galvo and Galvo-Resonant dual scanning modes. To address the femtosecond pulse dispersion issue, a pulse compression and dispersion compensation device based on a prism pair is designed, and a multi-pass cell (MPC) compression scheme is further analyzed, providing technical support for enhancing multiphoton excitation efficiency.

Finally, in terms of fluorescence lifetime data processing and system characterization, this thesis not only establishes the single-channel and three-channel synchronized time-correlated single-photon counting (TCSPC) acquisition pipelines and mitigates hardware dead-time limitations, but also develops a pixel-wise lifetime fitting algorithm with accuracy comparable to commercial software. Furthermore, to continuously improve the foundational support for multiphoton imaging, a two-photon interferometric autocorrelator based on a UV-responsive wide-bandgap semiconductor detector is independently built for the precise characterization of ultrafast pulses. A dual-ellipse cascaded equal-optical-path scanning architecture is also proposed, effectively correcting the field-of-view distortion and scan rate non-uniformity caused by traditional scanning configurations.

In summary, the work in this thesis completes a comprehensive construction ranging from cutting-edge scanning theory analysis, innovative opto-mechanical system design, multiphoton excitation, and pulse modulation, to fluorescence lifetime acquisition and custom data processing pipelines. It lays a solid hardware and software foundation for subsequent system optimization and application expansion oriented toward in vivo neural activity recording and multimodal functional imaging.
>>>>>>> c4bb8a8e4d07c80bffd10fe9aeb672bafe56fbee

\KEYWORDS{Multiphoton Microscopy, Multi-Color Confocal Microscopy, Fluorescence Lifetime Imaging, Time-Correlated Single-Photon Counting, Dispersion Compensation}% 英文关键词

\pagestyle{enfrontmatterstyle}%
\cleardoublepage\pagestyle{frontmatterstyle}%

%---------------------------------------------------------------------------%
